%!TEX root = main.tex
%=========================================================

\section{Emerging overlays}
\label{sec:adaptation}

We now present how overlays emerge (and \emph{dissolve}) autonomously in regions of a MANET where all nodes initially use CF.
We start by presenting how nodes leverage the MAC protocol to collect information about their environment.
An adaptation \emph{policy} triggers the creation of overlays based on thresholds over the collected \emph{observables}
The decision is collectively enforced for the necessary lifetime of the overlay, by piggybacking control information over existing message broadcasts.

\subsection{Observables: Density and Dynamism}

The emergence of an overlay in a MANET region (and its later dissolution) is driven by observed density and mobility.
Those could be measured explicitly by using periodic probing messages, as used by Tseng et~al.~\cite{2001:AAR:876878.879323} to adapt the parameters of controlled flooding.
The use of extra messages go against our objective of reducing the costs and risk of collisions.
Instead, we observe that the necessary information can be collected passively from the MAC protocol, BMW~\cite{tang2001mac}.
CTS requests are sent for every local broadcast.
Nodes in range reply with RTS frames, enabling the collection of their identities.
We collect the set of repliers for each RTS/CTS exchange, and keep a limited-time history of these neighborhoods (last 10 seconds in our implementation).

Density for a node $n_i$ is measured as its degree, or number of observed neighbors.
The measurements in the time window for $n_i$ form a sequence $d_{-x},d_{-w},d_{-v}\dots,d_0$, where $x\leq 10$ is the time of the latest measurement. 
Using only $d_0$ may lead to sudden fluctuations that do not reflect medium-term increases in density around the node.
We compute instead the density average over periods of time, e.g. $\bar{d}_{-1},\bar{d}_{-2},\bar{d}_{-5},\bar{d}_{-10}$.
As measurements are not taken with a fixed frequency, we weight each measurement using its validity window (e.g., $d_{-w}-d_{-x}$).

Mobility is measured indirectly through the \emph{stability} of $n_i$'s neighborhood, measured as the average time its neighbors are observed.
Again, we compute stability over varying periods of time, resulting in $\bar{s}_{-1},\bar{s}_{-2},\bar{s}_{-5},\bar{s}_{-10}$.

\subsection{Switching thresholds}

Nodes decide to emerge an overlay based on both $\bar{s}_{-\delta}$ and $\bar{d}_{-\delta}$ for a time window $\delta$.
In order to avoid oscillations where nodes would switch back and forth between using CF and MPR, we use a hysteresis, with two pairs of thresholds: 
	$({s}_{\text{min}},{d}_{\text{emerge}})$ 
and $({s}_{\text{min}},{d}_{\text{keep}})$, where
	${d}_{\text{keep}} < {d}_{\text{emerge}}$.
Criteria ${s}_{\text{min}}$ is a minimal stability required in all cases.
The two density criteria ${d}_{\text{emerge}}$ and ${d}_{\text{keep}}$ indicate when to respectively emerge and maintain an overlay.
The values of these criteria depend on the nature of the MANET, the wireless range of its nodes, and deployment factor.
We define appropriate values for our test scenarios in our evaluation.

\subsection{Switching policies}

A simple policy is to let nodes decide independently when to switch from CF to MPR, when the locally observed value for $(\bar{s}_{-\delta},\bar{d}_{-\delta})$ exceeds both thresholds $({s}_{\text{min}},{d}_{\text{emerge}})$.
As long as these values remain over thresholds $({s}_{\text{min}},{d}_{\text{keep}})$, the node continues using MPR.
While this \emph{single-node} policy will lead homogeneously dense and stable regions to switch to MPR eventually, it has the risk of creating small regions using MPR surrounded by CF zones, when nodes density is close but does not reach ${d}_{\text{emerge}}$.

%!TEX root = main.tex
%=========================================================

\begin{algorithm}[t]
\begin{footnotesize}
	\begin{algorithmic}[1]
	
	\BeginConstants
		\State $\Delta$: time \Comment{period of sending of \texttt{hello} messages for MPR}
		\State ${s}_{\text{min}},{d}_{\text{emerge}},{d}_{\text{keep}}$ \Comment{policy thresholds}
	\EndConstants
	
	\BeginVariables
		\State $\bar{s}_{-\delta}$, $\bar{d}_{-\delta}$ \Comment{observables maintained by MAC protocol}
		\State $P$ $\in$ \{ CF, MPR \} \Comment{currently running protocol}
		\State $t_\text{MPR}$ \Comment{when running MPR, dissolution timer}
	\EndVariables
	
	\smallskip
	
	\ProcedureAlt{emerge(): }
		\Return ($\bar{s}_{-\delta} \geq {s}_{\text{min}} \wedge \bar{d}_{-\delta} \geq {d}_{\text{emerge}}$)
		\label{algorithm:coordinated:line:emerge_condition}
	\EndProcedureAlt
	
	\ProcedureAlt{keep(): }
		\Return ($\bar{s}_{-\delta} \geq {s}_{\text{min}} \wedge \bar{d}_{-\delta} \geq {d}_{\text{keep}}$)
		\label{algorithm:coordinated:line:keep_condition}
	\EndProcedureAlt	
	
	\BeginUpon{processing of message $m$ for the first time}
		\If{emerge()} \Comment{act as a proposer}
			\State $m.$\texttt{em} $\gets$ \texttt{true}, useMPR()
		\ElsIf{$m.$\texttt{em}} \Comment{react to proposal:}
			\If{keep()} useMPR() \Comment{start or renew MPR use, or}
			\Else ~ $m.$\texttt{em} $\gets$ \texttt{false} \Comment{do not propagate proposal further}
			\EndIf
		\EndIf
	\EndUpon
	
	\ProcedureAlt{useMPR()}
		\If{$P =$ CF}
			\State MPR.init() \Comment{initialize \texttt{hello} messages}
			\State wait($2\times\Delta$), then $P \gets$ MPR \Comment{stabilize then switch}
		\EndIf
		\State $t_\text{MPR} \gets $ currentTime() $+ 10\times\Delta$ \Comment{set dissolution timer}
	\EndProcedureAlt
	
	\BeginTimerExpires{$t_\text{MPR}$}{\textbf{ or if no MPR hello received in} $2\Delta$}
		\State MPR.terminate()
		\State $P \gets$ CF \Comment{dissolve overlay and return to using CF}
	\EndTimerExpires

\end{algorithmic}
\end{footnotesize}
\caption{Coordinated adaptation on node $n_i$}\label{algorithm:coordinated}
\end{algorithm}



% \begin{algorithm}[t]
% \begin{footnotesize}
% 	\begin{algorithmic}[1]
%
% 	\BeginConstants
% 		\State $\Delta$: time \Comment{period of sending of \texttt{hello} messages for MPR}
% 		\State ${s}_{\text{min}},{d}_{\text{emerge}},{d}_{\text{keep}}$ \Comment{policy thresholds}
% 	\EndConstants
%
% 	\BeginVariables
% 		\State $\bar{s}_{-\delta}$, $\bar{d}_{-\delta}$ \Comment{observables maintained by MAC protocol}
% 		\State $P$ $\in$ \{ CF, MPR \} \Comment{currently running protocol}
% 		\State $t_\text{MPR}$ \Comment{when running MPR, dissolution timer}
% 	\EndVariables
%
% 	\smallskip
%
% 	\ProcedureAlt{emerge(): }
% 		\Return ($\bar{s}_{-\delta} \geq {s}_{\text{min}} \wedge \bar{d}_{-\delta} \geq {d}_{\text{emerge}}$)
% 		\label{algorithm:coordinated:line:emerge_condition}
% 	\EndProcedureAlt
%
% 	\ProcedureAlt{keep(): }
% 		\Return ($\bar{s}_{-\delta} \geq {s}_{\text{min}} \wedge \bar{d}_{-\delta} \geq {d}_{\text{keep}}$)
% 		\label{algorithm:coordinated:line:keep_condition}
% 	\EndProcedureAlt
%
% 	\BeginUpon{processing of message $m$ \emph{for the first time}}
% 		\If{emerge()}
% 			\State $m.$\texttt{em} $\gets$ \texttt{true}
% 			\If{$P =$ CF} switchToMPR() \EndIf
% 		\ElsIf{$m.$\texttt{em}}
% 			\If{keep()}
% 				\If{$P =$ CF} switchToMPR()
% 				\Else ~$t_\text{MPR} \gets $ currentTime() $+ 10\times\Delta$
% 				\EndIf
% 			\Else
% 				~ $m.$\texttt{em} $\gets$ \texttt{false}
% 			\EndIf
% 		\EndIf
% 	\EndUpon
%
% 	\ProcedureAlt{switchToMPR()}
% 		\State MPR.init() \Comment{initialize the sending of \texttt{hello} messages by MPR}
% 		\State wait($2\times\Delta$), then $P \gets$ MPR \Comment{stabilize and switch}
% 		\State $t_\text{MPR} \gets $ currentTime() $+ 10\times\Delta$ \Comment{set dissolution timer}
% 	\EndProcedureAlt
%
% 	\BeginTimerExpires{$t_\text{MPR}$}{\textbf{ or if no MPR hello received in} $2\Delta$}
% 		\State MPR.terminate()
% 		\State $P \gets$ CF \Comment{dissolve overlay and return to using CF}
% 	\EndTimerExpires
%
% \end{algorithmic}
% \end{footnotesize}
% \caption{Coordinated adaptation on node $n_i$}\label{algorithm:coordinated}
% \end{algorithm}


Our second \emph{coordinated} policy implements lightweight coordination between nodes, by piggybacking a single bit of information in the header of existing messages broadcast.
This bit, named \texttt{em} for \emph{emerge!}, carries the information that 



To switch from CF to MPR, nodes nee





The decision to emerge an overlay in a region of the MANET is based primarily on the organization of nodes in space (the  density in the neighborhood of a node) and on their dynamism (i.e., witnessed as frequent reconfigurations in the neighborhood of a node).
The collection of up-to-date information would require constant probing by each node of its surrounding, which would defeat the purpose of emergent overlays of reducing costs and improving reliability through less collisions.
Instead, we adopt a passive approach using existing communications.
We collect the following \emph{observables} about the surroundings of a node.
Each of them is averaged over a configurable number of observations in order to privilege long-term changes in the deployment environment.
\begin{itemize}
	\item The density is obtained from the MAC protocol BMW. It is the number of nodes who reply to local broadcast RTS requests.
	\item 
\end{itemize}


observables:
\begin{itemize}
	\item input source: MAC protocol, responses to RTS frames + only when MPR protocol is running, also responses to direct control frames
	\item used to derive local:
	\begin{itemize}
		\item density (how many nodes in range)
		\item clustering between one-hop neighbors
		\item stability (use a figure: we probe the network regularly through existing mac-layer messages, we can keep snapshots of the neighbors, and estimate how much time we see a neighbor)
	\end{itemize}
\end{itemize}

policies:
\begin{enumerate}
	\item switching criteria based on nodes degree: collective decision where nodes with similar values of degree chose the same protocol. Currently, the policy works as follows:
	\begin{itemize}
		\item 3 thresholds (low, medium, high) for local clustering coefficient $c$, for: 
		$c \in [0, \frac{1}{3}) \rightarrow t_1 = 0$ (implies controlled flooding), 
		$c \in [\frac{1}{3}, \frac{2}{3}) \rightarrow t_2 = 1 $ (these are usually border nodes), 
		$c \in [\frac{2}{3}, 1) \rightarrow t_3 = 3 $ (implies MPR: overlay-based protocol).
		\item remind that nodes know their degree using the clustering observable. Switching procedure is as follows:
		\begin{enumerate}
			\item the source of a message $m$ attach a frame $f$ that contains $t_i \in [0, 1, 2]$ and start a timer $s$ to switch of protocol
			\item receivers of $m$ detect $f$ and make the difference $d$ between the threshold in $m$ and their local threshold; receivers start a timer $s$ to switch of protocol f.$t_i$ only if $d$ is equal to the threshold of the source node, otherwise:
			\begin{itemize}
				\item replace f.$t_i$ with $n_i$.$t_j$ (local threshold) in $m$ and forwards $m$
				\item start timer $s$ if needed
			\end{itemize}
			\item when $s$ expires all nodes switch of protocol
		\end{enumerate}
	\end{itemize}
\end{enumerate}

\er{
The goal of this section is to present how adaptation is performed.
It should start with a general discussion of the control loop: local sensing, collection of information, periodic checking of policies, trigger of adaptation actions (here, emerging an overlay). If we do it, explain that these phases may happen on a single node (single-node adaptation) or through the collaboration of several nodes in a neighborhood (multi-node adaptation).
Suggested outline:
\begin{itemize}
	\item Adaptation control loop: Define sensing, information collection, policies checking, adaptation actions.
	\item Define single-node adaptation and multi-node adaptation. Do not give details on how the latter is done, say this is described in a later subsection.
	\item Subsection on sensing: Explain that we are limited in the information we can collect as we do not want to send probes to the neighbors, which would consume energy, increase the amount of collisions, and defeat the purpose. Henceforth we rely on what we can observe through existing interactions. In other word, we adopt passive sensing.
	\item Define how we can classify the information we collect, and what is their nature: snapshot (instantaneous information about the surrounding of the node, at a given time), or continuous (history information about the evolution of some observable, over a window of time or a number of measurements), relative (from 0 to 1) or absolute, complete (e.g. the number of nodes in the neighborhood) or partial, and other criteria as necessary/possible.
	\item Present a table with the set of observables we will collect and their classification according to the taxonomy presented in the previous point. We want at least 4-5 observables in order to form a diverse set of policies.
	\item Add a discussion about the difference in observables depending on the protocol running, which makes our adaptation problem special: the policies may get different information for switching from one protocol to the other, than they get for doing the opposite. (keep if this is justified in practice for the policies we do present)
	\item The next subsection is devoted to the design of a series of adaptation policies, from very simple, naive ones using a single observable to more advanced ones using multiple observables. They have to be presented using a table, with conjunction of rules based on thresholds, trigger time (periodic, when there is a new observable value obtained, etc.), and nature (single vs. multi-node)
	\item Discuss each policy and give the design principle, the objective, rationale and expected benefits and limitations.
	\item The multi-node policies should use figures to ease their explanation, and there can also be an algorithm showing how the nodes interact to collect information, and reach a decision.
\end{itemize}
}

\er{
Suggestion to discuss:
it would also be possible to have two parts, one where we only discuss single-node policies, and a second one where we build an extension for multi-node policies.
This could simplify the presentation.
The multi-node policy requires communication.
We need to assess whether this communication can happen on top of existing messages or not (would be ideal) or if it must use its own control messages.
}

% NOTE: These papers may help in the construction of a policy about density and mobility:
%\cite{closure_coef}: local closure coefficient used as metric of density in our policies of adaptation
% to check rapidly, as this is a C-rank journal: \url{https://link.springer.com/article/10.1007/s11277-016-3213-0} proposes to have two modes of operation depending on the mobility of the nodes (so quite linked to us)
%to check as a possible example of a parameters adaptation (to group with the others of the same kind): \url{https://ieeexplore.ieee.org/abstract/document/7344460}
